% ============================================================
% 英文简历模板（resume-en，现代单栏风）
% ------------------------------------------------------------
% 样式由模板自绘（无第三方简历类的依赖），编译引擎：pdflatex。
% 修改指南：
%   1. "Header"区填写姓名与联系方式（图标行按需增删）；
%   2. 各分区直接复制 \entry 与 itemize 条目修改内容；
%   3. 主题色在 \definecolor 处修改，全局生效；
%   4. 目标一页：删减内容而不是缩小字号；
%   5. 中文求职场景请改用 resume-cn 模板。
% ============================================================
\documentclass[10pt, a4paper]{article}

% ---- 字体与页面 ----
\usepackage[T1]{fontenc}
\usepackage[default]{sourcesanspro}       % 现代无衬线字体（TeX Live 自带）
\usepackage[
  top=1.4cm, bottom=1.4cm, left=1.6cm, right=1.6cm
]{geometry}
\pagestyle{empty}                          % 简历不显示页码
\setlength{\parindent}{0pt}

% ---- 依赖宏包 ----
\usepackage{xcolor}
\usepackage{enumitem}
\usepackage{titlesec}
\usepackage{fontawesome5}                  % 联系方式图标
\usepackage{hyperref}

% 主题色：换色改这里（如深蓝 / 墨绿 / 深红）
\definecolor{primary}{RGB}{20, 60, 110}
\definecolor{graytext}{RGB}{100, 100, 100}

\hypersetup{colorlinks=true, urlcolor=primary, linkcolor=primary}

% ---- 分区标题：加粗着色 + 下划细线 ----
\titleformat{\section}
  {\large\bfseries\color{primary}}{}{0em}{}
  [\vspace{-0.5em}{\color{primary!35}\rule{\textwidth}{0.8pt}}]
\titlespacing{\section}{0pt}{10pt}{6pt}

% ---- 列表样式：紧凑、圆点、左缩进小 ----
\setlist[itemize]{
  nosep,
  leftmargin=1.2em,
  label=\textcolor{primary}{\small$\bullet$},
}

% ---- 经历条目命令：\entry{职位/学位}{单位/学校}{时间}{地点} ----
\newcommand{\entry}[4]{%
  \textbf{#1} \textcolor{graytext}{$\vert$} #2
  \hfill \textcolor{graytext}{#3 \, #4}\par
  \vspace{0.15em}
}

% ---- 技能条目命令：\skill{类别}{内容} ----
\newcommand{\skill}[2]{\textbf{#1:} #2\par}

\begin{document}

% ==================== Header ====================
\begin{center}
  {\Huge\bfseries\color{primary} Your Name}\\[0.6em]
  % 联系方式：按需保留/删除图标项，\faXxx 为 fontawesome5 图标
  \faEnvelope\ \href{mailto:you@example.com}{you@example.com} \quad
  \faPhone\ +86 138-0000-0000 \quad
  \faGithub\ \href{https://github.com/yourname}{github.com/yourname} \quad
  \faLinkedin\ \href{https://linkedin.com/in/yourname}{linkedin.com/in/yourname}
\end{center}
\vspace{-0.4em}

% ==================== Summary ====================
\section{Summary}
% 两三行概括定位：专业方向 + 核心技能 + 亮点成果
Computer science senior specializing in machine learning systems.
Experienced in building data pipelines and productionizing models;
author of 2 conference papers.

% ==================== Education ====================
\section{Education}
\entry{M.S. in Computer Science}{Example University}{2024.09 -- 2027.06 (expected)}{Beijing}
\begin{itemize}
  \item GPA: 3.8/4.0; focus: machine learning and computer vision.
\end{itemize}
\entry{B.S. in Computer Science}{Example University}{2020.09 -- 2024.06}{Beijing}
\begin{itemize}
  \item Ranked top 5\%; received the National Scholarship (2022).
\end{itemize}

% ==================== Experience ====================
\section{Experience}
\entry{Machine Learning Intern}{Tech Company, Beijing}{2025.06 -- 2025.09}{Beijing}
\begin{itemize}
  \item Built a data cleaning pipeline processing 10M+ records daily.
  \item Improved model inference speed by 30\% through quantization.
\end{itemize}
\entry{Research Assistant}{AI Lab, Example University}{2023.10 -- 2024.06}{Beijing}
\begin{itemize}
  \item Designed experiments for the paper accepted at ExampleConf 2024.
\end{itemize}

% ==================== Projects ====================
\section{Projects}
\entry{LaTeXcellence}{AI-powered LaTeX assistant}{2026}{Personal project}
\begin{itemize}
  \item Developed a FastAPI + React tool for AI-assisted proofreading of \LaTeX{} documents.
  \item Repo: \href{https://github.com/yourname/project}{github.com/yourname/project}.
\end{itemize}

% ==================== Skills ====================
\section{Skills}
\skill{Languages}{Python, C++, TypeScript, SQL}
\skill{Frameworks}{PyTorch, FastAPI, React, Docker}
\skill{Languages (spoken)}{English (fluent, TOEFL 105), Chinese (native)}

% ==================== Honors ====================
\section{Honors \& Awards}
\begin{itemize}
  \item National Scholarship, Ministry of Education, 2022.
  \item Gold Medal, Regional Programming Contest (ICPC), 2023.
\end{itemize}

\end{document}
